\documentclass[12pt]{article}
\usepackage[T1]{fontenc}
\usepackage[utf8]{inputenc}
\usepackage{lmodern}
\usepackage{textcomp}
\usepackage{enumitem}
\usepackage{geometry}
\usepackage{graphicx}
\usepackage{amsmath, amssymb}
\geometry{margin=1in}
\begin{document}
\begin{center}
Constitutional Law - Undergraduate Level Assessment 19 \
Total Marks: 40 \
Answer all questions.
\end{center}

\section*{Multiple Choice (10 marks)}
\begin{enumerate}[label=\arabic*.]
\item Why does Parfit oppose equality?
\begin{enumerate}[label=(\alph*)]
    \item He argues than an unequal society is inevitable.
    \item He claims that by giving priority to the needs of the poor, we can increase equality.
    \item He asserts that we are each responsible for our poverty.
    \item He rejects the idea of equality altogether.
\end{enumerate}
\item What is an act jure imperii?
\begin{enumerate}[label=(\alph*)]
    \item An act is jure imperii when undertaken by an international organisation
    \item An act is jure imperii when undertaken in an official State capacity
    \item All acts undertaken by State officials are acts jure imperii
    \item An act is jure imperii when undertaken by a State corporation
\end{enumerate}
\item Dworkin argues that it is only a conception of equality of resources that can secure the ideal of equality of welfare. How does he suggest this aspect of equality to be measured?
\begin{enumerate}[label=(\alph*)]
    \item When no-one would prefer another's bundle of resources to his or her own.
    \item By reference to the ownership of private property.
    \item By the amount of income tax paid by individuals.
    \item When the community determines that equality has been achieved.
\end{enumerate}
\item What is the purpose of politics, according to Aristotle?
\begin{enumerate}[label=(\alph*)]
    \item To advance the interests of politicians.
    \item To produce virtuous citizens and encourage righteousness in individuals.
    \item To obtain power.
    \item To secure justice by abolishing slavery.
\end{enumerate}
\item Was the use of armed force permitted prior to the United Nations Charter?
\begin{enumerate}[label=(\alph*)]
    \item Armed force was prohibited
    \item Armed force was permitted with no restrictions
    \item Armed force was permitted subject to few restrictions
    \item Armed force was not regulated under international law prior to 1945
\end{enumerate}
\end{enumerate}

\section*{True / False (10 marks)}
\textit{Answer True or False}
\begin{enumerate}[label=\arabic*.]
\setcounter{enumi}{5}
\item True or False: The idea of 'natural law' first appeared in Greek thinking, influencing later legal and philosophical developments.
\item True or False: Friedrich Kelsen was the first jurist to study the comparative aspect of law, focusing on legal positivism.
\item True or False: Dworkin's notion of law as integrity does not support the idea that it opens the door to authoritarianism.
\item True or False: The Analytical School posits that a custom becomes law when it is recognized and endorsed by judicial decisions.
\item True or False: Codification of law simplifies legal interpretation and makes it more accessible to the public.
\end{enumerate}

\section*{Long Form Response (20 marks)}
\textit{Answer each question with a clear and complete explanation.}
\begin{enumerate}[label=\arabic*.]
\setcounter{enumi}{10}
\item Discuss the issue of 'fragmentation' in international law, focusing on how the coexistence of different legal regimes can lead to divergent rules and implications for governance.
\item Discuss the concept of civil law, focusing on how it establishes rights between individuals and the mechanisms it provides for redress in cases of rights violations.
\end{enumerate}

\vfill
\noindent\textit{End of Paper}
\end{document}